\subsection{Review of Finance (Cont.)}

\content{
        Formulas so far: \\
        \fbox{\begin{minipage}{\textwidth}
        For these formulas, we have the following: 
        $P_0$ is the principal. 
        $r$ is the interest rate. 
        $t$ is the number of time intervals (specified by problem) which
        have passed. 
        $A$ is the amount paid back.
        $I$ is the interest.

        \begin{center}
        \begin{tabular}{|c|c|}
        \hline
        \T\B Simple Interest &  $ A = P_0(1+r) $ \\
        \hline
        \T\B Simple Interest over Time & $  A = P_0(1+rt) $ \\
        \hline
        \end{tabular}
        \end{center}
        \end{minipage}}
        \fbox{\begin{minipage}{\textwidth}
        For these formulas, we have the following: 
        $P_0$ is the principal, or amount of a loan. 
        $r$ is the annual interest rate.
        $t$ is the number of years which have passed. 
        $A$ is the total amount.
        $k$ is the number of times interest is compounded per year. 
        $d$ is the amount we deposit (or withdraw) from the account regularly,
        or the payment amount on a loan.

        \begin{center}
        \begin{tabular}{|c|c|}
        \hline
        \T\B Compound Interest & $ A =
        P_0\left(1+\frac{r}{k}\right)^{kt} $ \\
        \hline
        \T\B Savings Annuity & $ A =
        d\frac{\left(\left(1+\frac{r}{k}\right)^{kt}-1\right)}{\frac{r}{k}} $\\
        \hline
        \T\B Payout Annuity/Loans & $P_0 =
        d\frac{\left(1-\left(1+\frac{r}{k}\right)^{-kt}\right)}{\frac{r}{k}}$ \\
        \hline
        \end{tabular}
        \end{center}
        \end{minipage}}
}{% presentation
}{% nopresentation
}{% comments
}

\content{
        \begin{example}
        Marcy received an inheritance of \$20,000, and invested it at 6\% interest. She is
        going to use it for college, withdrawing money for tuition and expenses each
        quarter. How much can she take out each quarter if she has 3 years of school
        left?        
        \end{example}
}{% presentation
\vspace{2in}
}{% nopresentation
\vspace{2in}
}{% comments
}

\content{
        \begin{example}
        Paul wants to buy a new car. Rather than take out a loan, he decides to save
        \$200 a month in an account earning 3\% interest compounded monthly. How
        much will he have saved up after 3 years?        
        \end{example}
}{% presentation
\vspace{2in}
}{% nopresentation
\vspace{2in}
}{% comments
}

\content{
        \begin{example}
        Keisha is managing investments for a non-profit company. They want to invest
        some money in an account earning 5\% interest compounded annually with the
        goal to have \$30,000 in the account in 6 years. How much should Keisha
        deposit into the account?        
        \end{example}
}{% presentation
\vspace{2in}
}{% nopresentation
\vspace{2in}
}{% comments
}

\content{
        \begin{example}
        Miao is going to finance new office equipment at a 2\% rate over a 4 year term.
        If she can afford monthly payments of \$100, how much new equipment can she
        buy?        
        \end{example}
}{% presentation
\vspace{2in}
}{% nopresentation
\vspace{2in}
}{% comments
}

\content{
        \begin{example}
        How much would you need to save every month in an account earning 4\%
        interest to have \$5,000 saved up in two years?        
        \end{example}
}{% presentation
\vspace{2in}
}{% nopresentation
\vspace{2in}
}{% comments
}

\content{
        \begin{example}
        A friend asks for a small loan of \$100 and agrees to pay it back with
        10\% interest in a week. For each week that goes by without payment,
        another \$10 is added to what they owe you. Assuming that they don't pay
        you back for four weeks, how much will they owe you? 
        \end{example}
}{% presentation
\vspace{2in}
}{% nopresentation
\vspace{2in}
}{% comments
}
